\section*{Bab \ref{ch:b1:phase-changes} --- Perubahan Fase Zat Murni}

\begin{solution}{exo:b1:phase-changes:1}
(a) Uap: \qty{0.5}{kPa} berada di bawah $P_s(\qty{20}{\degreeCelsius}) = \qty{2.3}{kPa}$.
(b) Padat. (c) Di atas $T_c = \qty{374}{\degreeCelsius}$: fluida superkritis, tanpa
pembedaan cair--gas. (d) \omterm{prop:b1:phase-changes:PTdiagram}{Titik tripelnya}: ketiganya. Genangannya
menguap dari permukaannya selama tekanan uap di udaranya
di bawah $P_s(\qty{20}{\degreeCelsius})$; pendidihan hanyalah kasus ketika
gelembung dapat terbentuk di dalam cairannya.
\end{solution}

\begin{solution}{exo:b1:phase-changes:2}
$0.1 \times 2257 + 0.1 \times 4.18\,(100 - T_f) = 1.0 \times 4.18\,(T_f - 20)$:
$351.1 = 4.60\,T_f$, $T_f = \qty{76}{\degreeCelsius}$. Dengan air panas
sebagai gantinya: $100 - T_f = 10(T_f - 20)$, $T_f = \qty{27}{\degreeCelsius}$ ---
\omterm{def:b1:phase-changes:phase}{kalor latennya} setara dengan \qty{540}{K} kalor sensibel.
\end{solution}

\begin{solution}{exo:b1:phase-changes:3}
$v = \qty{2.0e-3}{m^3/kg}$; $x = (2.0 - 1.09) \times 10^{-3}/0.392 = 2.3 \times 10^{-3}$:
\qty{2.3}{g} uap yang mengisi \qty{0.91}{L}; \qty{0.998}{kg} cairannya
menempati $0.998 \times 1.09 = \qty{1.09}{L}$.
\end{solution}

\begin{solution}{exo:b1:phase-changes:4}
$2257/373 = \qty{6.05}{kJ/(kg.K)}$, yakni $\times 0.018 = \qty{109}{J/(mol.K)}$;
peleburannya $334/273 \times 0.018 = \qty{22}{J/(mol.K)}$. Nilai $109$ milik air jauh
di atas $88$ milik Trouton: ikatan hidrogen membuat air cair lebih teratur
daripada cairan biasa, sehingga menguapkannya memperoleh lebih banyak ketakteraturan.
\end{solution}

\begin{solution}{exo:b1:phase-changes:5}
$1/T = 1/373.15 - \ln(P/P_0)/4890$. Everest: $\ln(0.33/1.013) = -1.12$,
$1/T = 2.680 \times 10^{-3} + 2.29 \times 10^{-4}$, $T = \qty{344}{K} =
\qty{71}{\degreeCelsius}$. Pada \qty{0.80}{bar}: $\ln = -0.236$, $T = \qty{366}{K} =
\qty{93}{\degreeCelsius}$. Pada \qty{71}{\degreeCelsius} putih telurnya mengeras perlahan
dan kuningnya nyaris tidak: telur \qty{3}{min} menjadi telur
\qty{10}{min} atau lebih.
\end{solution}

\begin{solution}{exo:b1:phase-changes:6}
$\ln(1.8/1.013) = 0.575$: $1/T = 2.680 \times 10^{-3} - 1.18 \times 10^{-4}$,
$T = \qty{390}{K} = \qty{117}{\degreeCelsius}$; pada \qty{2.5}{bar}: $\ln = 0.903$,
$T = \qty{401}{K} = \qty{128}{\degreeCelsius}$. Uapnya: $\qty{2}{L}/0.98 =
\qty{2.0}{g}$ --- pancinya nyaris tak memuat apa pun selain air dan sedikit
uap.
\end{solution}

\begin{solution}{exo:b1:phase-changes:7}
$\dd P/\dd T = 3.34 \times 10^5/[273 \times (-9.1 \times 10^{-5})] = \qty{-1.3e7}{Pa/K}
= \qty{-130}{bar/K}$: es mencair pada \qty{-1}{\degreeCelsius} di bawah kira-kira
\qty{130}{bar}. Di bawah \qty{3}{km} es, $P = 917 \times 9.8 \times 3000 =
\qty{270}{bar}$: leburnya pada \qty{-2}{\degreeCelsius}; dasar sebuah gletser dapat berada pada
titik leburnya sementara permukaannya jauh lebih dingin, dan selaput
air lelehnya membuat lapisan esnya meluncur.
\end{solution}

\begin{solution}{exo:b1:phase-changes:8}
$W = -P(v_v - v_l) = -1.013 \times 10^5 \times 1.672 = \qty{-169}{kJ}$ (airnya
mendorong mundur atmosfernya); $Q = L_v = \qty{2257}{kJ}$; $\Delta U = 2257 - 169
= \qty{2088}{kJ}$; $\Delta H = \qty{2257}{kJ}$; $\Delta S = 2257/373 = \qty{6.05}{kJ/K}$.
Dari nyala api pada \qty{1500}{K}: $S_{\mathrm{tukar}} = 2257/1500 = \qty{1.50}{kJ/K}$,
yang tercipta \qty{4.5}{kJ/K} --- sebagian besar mutu nyala apinya terbuang untuk memanaskan
air.
\end{solution}

\begin{solution}{exo:b1:phase-changes:9}
$P_s(\qty{22}{\degreeCelsius}) \approx 2.34 + 0.4 \times 0.83 = \qty{2.67}{kPa}$,
$P_v = 0.70 \times 2.67 = \qty{1.87}{kPa}$; \omterm{ex:b1:phase-changes:dew}{titik embunnya} tempat $P_s = 1.87$: $15 +
5 \times 0.17/0.64 = \qty{16}{\degreeCelsius}$. Massanya: $\rho_v = 1870 \times 0.018/
(8.314 \times 295) = \qty{13.7}{g/m^3}$, \qty{0.69}{kg} di dalam ruangannya. Kacanya pada
\qty{12}{\degreeCelsius}: $P_s \approx \qty{1.42}{kPa} < 1.87$, ia berembun. Pada
\qty{25}{\degreeCelsius} kelembapannya turun ke $1.87/3.17 = 59\%$.
\end{solution}

\begin{solution}{exo:b1:phase-changes:10}
(a) $h$ adalah fungsi keadaan: menempuh padat $\to$ cair $\to$ uap di
\omterm{prop:b1:phase-changes:PTdiagram}{titik tripelnya} atau padat $\to$ uap secara langsung memberikan $\Delta h$ yang sama:
$L_s = L_f + L_v$. (b) Dengan $v_s, v_l \ll v_v$ kedua kemiringannya $L/(Tv_v)$ pada
$T$ dan $v_v$ yang sama: nisbahnya $L_s/L_v > 1$. (c) $v_v = RT/(MP) =
8.314 \times 273.16/(0.018 \times 611) = \qty{206}{m^3/kg}$; penguapannya
$2.50 \times 10^6/(273.16 \times 206) = \qty{44}{Pa/K}$, sublimasinya $2.83 \times 10^6/
(273.16 \times 206) = \qty{50}{Pa/K}$. (d) Di bawah \qty{0}{\degreeCelsius} pada
\qty{1}{bar} tak ada cairan; bila tekanan uap udaranya di bawah
$P_s$ di atas es (\qty{260}{Pa} pada \qty{-10}{\degreeCelsius}) saljunya
menyublim --- angin dingin yang kering memakannya.
\end{solution}

\begin{solution}{exo:b1:phase-changes:11}
(a) $P_s(\qty{105}{\degreeCelsius}) \approx 1.013 + 5 \times 0.036 = \qty{1.19}{bar}$:
$\Delta P = \qty{0.18}{bar}$, $r^\ast = 2 \times 0.059/(1.8 \times 10^4) = \qty{6.6}{\micro m}$.
(b) $\Delta P = \qty{3.6}{kPa}$, $r^\ast = \qty{33}{\micro m}$. (c) $x = c\,\Delta T/L_v =
4.18 \times 5/2257 = 0.93\%$: \qty{9.3}{g} tiap liter, yakni $9.3 \times 10^{-3}
\times 1.67 = \qty{16}{L}$ uap yang dihasilkan dalam sepersekian sekon di dalam
sebuah cangkir --- airnya tersembur keluar sambil mendidih. (d) Batu didih menahan udara di dalam
pori-porinya: gelembung berukuran puluhan mikrometer sudah ada, sehingga pendidihannya bermula
begitu $P_s$ mencapai $P_0$.
\end{solution}

\begin{solution}{exo:b1:phase-changes:12}
(a) Setelah katupnya $h = \qty{256}{kJ/kg}$ pada \qty{-10}{\degreeCelsius}: $x =
(256 - 187)/(392 - 187) = 0.34$. (b) Evaporatornya $q_c = 392 - 256 =
\qty{136}{kJ/kg}$; kompresornya $w = 430 - 392 = \qty{38}{kJ/kg}$; kondensornya
$q_h = 256 - 430 = \qty{-174}{kJ/kg}$; $136 + 38 - 174 = 0$. (c) $e_p =
174/38 = 4.6$, $e_f = 136/38 = 3.6$; Carnot $313/50 = 6.3$ dan $263/50 =
5.3$: kira-kira $0.7$ dari yang ideal. (d) $\dot m = 6000/174000 = \qty{34}{g/s}$,
\qty{2.1}{kg} tiap menit.
\end{solution}

\begin{solution}{pb:b1:phase-changes:1}
\textbf{1.} Pada \qty{1}{atm} air mendidih pada \qty{100}{\degreeCelsius}; selama
cairannya masih ada, kalor yang dipasok menghasilkan uap pada suhu itu
(\omterm{def:b1:phase-changes:phase}{kalor laten}) alih-alih menghangatkan cairannya: suhunya terpasak
oleh tekanannya.

\textbf{2.} $\dd P_s/\dd T = L_v/[T(v_v - v_l)]$ dengan $v_l \ll v_v = RT/(MP_s)$
(uapnya \omterm{def:b1:kinetic-theory:model}{gas ideal}): $\dd P_s/P_s = (ML_v/R)\,\dd T/T^2$; dengan $L_v$ dianggap
tetap, integralkan: $\ln(P_s/P_0) = (ML_v/R)(1/T_0 - 1/T)$.

\textbf{3.} $1/T = 1/373.15 - \ln(1.8/1.013)/4890 = 2.562 \times 10^{-3}$: $T =
\qty{390}{K} = \qty{117}{\degreeCelsius}$.

\textbf{4.} $m = \Delta P\,A/g = 0.8 \times 10^5 \times 2.0 \times 10^{-5}/9.8 =
\qty{0.16}{kg}$.

\textbf{5.} $Q = 2.0 \times 4.18 \times 97 = \qty{811}{kJ}$; $t = 811/2.0 =
\qty{405}{s} \approx \qty{7}{min}$.

\textbf{6.} $2000/(2.2 \times 10^6) = \qty{0.91}{g/s} = \qty{55}{g/min}$ --- dua
liter habis dalam \qty{37}{min}. Kecilkan tungkunya sampai sekadar cukup untuk
membuat katupnya tetap mendesis: suhunya ditetapkan oleh tekanannya, bukan oleh
apinya.

\textbf{7.} $+\qty{17}{K}$: $2^{1.7} = 3.2$ kali lebih cepat, \qty{30}{min}
$\to$ \qty{9}{min}. Everest, $-\qty{30}{K}$: $2^{-3}$, empat jam.

\textbf{8.} \omterm{ex:b1:phase-changes:dew}{Kelembapan nisbi} $= P_v/P_s(T)$. $P_v = 0.60 \times 3.17 =
\qty{1.90}{kPa}$; $\rho_v = P_vM/(RT) = 1900 \times 0.018/(8.314 \times 298) =
\qty{13.8}{g/m^3}$.

\textbf{9.} $P_s(T_d) = \qty{1.90}{kPa}$: di antara \qty{15}{\degreeCelsius}
(\qty{1.70}) dan \qty{20}{\degreeCelsius} (\qty{2.34}): $T_d = 15 + 5 \times 0.20/0.64
= \qty{16.6}{\degreeCelsius}$. Rumput yang memancar ke langit malam, atau
botol dari \omterm{def:b1:heat-engines:machine}{lemari es}, turun di bawah $T_d$: uap yang bersentuhan
dengannya lalu berada di atas kejenuhan dan mengembun.

\textbf{10.} $TP^{(1-\gamma)/\gamma} = \text{tetap}$: $\dd T/T = \frac{\gamma - 1}{\gamma}
\dd P/P$; $\dd P = -\rho g\,\dd z$ dengan $\rho = PM/(RT)$: $\dd T/T = -\frac{\gamma -
1}{\gamma}\frac{Mg}{RT}\dd z$, sehingga $\dd T/\dd z = -\frac{\gamma - 1}{\gamma}\frac{Mg}{R} =
-g/c_p$ karena $c_p = \gamma R/[(\gamma - 1)M]$. Nilainya $9.8/1000 = \qty{9.8}{K/km}$.

\textbf{11.} $T - T_d$ menutup dengan $9.8 - 2 = \qty{7.8}{K/km}$: $z = (25 - 16.6)/
7.8 = \qty{1.1}{km}$.

\textbf{12.} $1 \times 2.45 = \qty{2.45}{kJ}$ tiap kilogram udara: $\Delta T =
\qty{2.45}{K}$. Pelepasannya mengimbangi sebagian pendinginan adiabatiknya: bongkahan
yang jenuh mendingin hanya \qtyrange{5}{6}{K/km}, tinggal lebih hangat dan
lebih ringan daripada sekelilingnya lalu terus naik --- \omterm{def:b1:phase-changes:phase}{kalor latennya} adalah
bahan bakar badai guruhnya.

\textbf{13.} $10^9 \times \qty{1}{g} = \qty{1000}{t}$; $10^6 \times 2.45 \times 10^6 =
\qty{2.5e12}{J} = \qty{0.68}{GW h}$ --- empat puluh menit sebuah pembangkit besar.

\textbf{14.} $P_v = 0.80 \times 3.17 = \qty{2.54}{kPa}$; $P_s(\qty{18}{\degreeCelsius})
\approx 1.70 + 0.6 \times 0.64 = \qty{2.08}{kPa} < 2.54$: ia berembun, sampai
kelembapannya turun di bawah $2.08/3.17 = 66\%$.

\textbf{15.} $P = P_s(\qty{20}{\degreeCelsius}) = \qty{2.3}{kPa}$; $v = \qty{2.0e-3}{m^3/kg}$;
$v_v = RT/(MP_s) = 8.314 \times 293/(0.018 \times 2340) = \qty{58}{m^3/kg}$; $x =
(2.0 - 1.0) \times 10^{-3}/58 = 1.7 \times 10^{-5}$: \qty{9}{mg} uap;
cairannya $0.50 \times 10^{-3} = \qty{0.50}{L}$, separuh bejananya.

\textbf{16.} $v = \qty{2.0e-3}{m^3/kg} < v_c$: isokornya menemui cabang
cairan kubahnya --- cairannya memuai lebih cepat daripada ia menguap, arasnya
\emph{naik}, dan bejananya penuh cairan ketika $v_l(T) =
\qty{2.0e-3}{m^3/kg}$, sedikit di atas \qty{360}{\degreeCelsius} (kira-kira
\qty{190}{bar}). Di luar itu, cairan yang hampir tak termampatkan dipanaskan pada volume
tetap: tekanannya memanjat puluhan bar tiap \omterm{def:b1:kinetic-theory:temperature}{kelvin} --- bejananya
pecah. Jangan pernah menyegel bejana yang penuh cairan.

\textbf{17.} $v = \qty{1.0e-2}{m^3/kg} > v_c$: arasnya \emph{turun}; tetes
terakhirnya pergi ketika $v_v(T) = 0.010$, di antara \qty{340}{\degreeCelsius} ($0.0108$)
dan \qty{350}{\degreeCelsius} ($0.0088$): kira-kira \qty{344}{\degreeCelsius},
\qty{150}{bar}.

\textbf{18.} $v = \qty{3.2e-3}{m^3/kg} \approx v_c$: meniskusnya tinggal dekat
setengah tingginya, menjadi makin rata dan makin pudar ketika kedua massa jenisnya menyatu,
fluidanya menjadi seperti susu (opalesensi kritis), dan pada \qty{374}{\degreeCelsius},
\qty{221}{bar} meniskusnya lenyap: satu \omterm{def:b1:phase-changes:phase}{fase}.

\textbf{19.} Tiga garis tegak: $v = 2.0 \times 10^{-3}$ meninggalkan kubahnya
melalui cabang cairannya, di kiri C; $10^{-2}$ melalui cabang uapnya,
di kanan C; $3.2 \times 10^{-3}$ melalui puncaknya.

\textbf{20.} Udaranya: \qty{0.5}{L} pada \qty{1.0}{bar}, \qty{293}{K}, dengan volume yang sama pada
\qty{373}{K}: $P_{\text{udara}} = 1.0 \times 373/293 = \qty{1.27}{bar}$; ditambah
$P_s = \qty{1.01}{bar}$: \qty{2.3}{bar} di dalamnya, \qty{1.3}{bar} di atas tekanan
di luarnya: $1.3 \times 10^5 \times 10^{-2} = \qty{1.3}{kN}$ pada tutupnya --- setara
berat \qty{130}{kg}.

\textbf{21.} \omterm{def:b1:phase-changes:metastable}{Lewat dingin}, metastabil: tak ada tapak \omterm{def:b1:phase-changes:metastable}{nukleasi} (air yang bersih,
tanpa debu, tanpa gerakan) bagi kristal yang pertama. Sebuah ketukan menumbuhkan inti es, yang
tumbuh menembus botolnya sambil melepaskan $L_f$; suhunya naik ke
\qty{0}{\degreeCelsius} lalu berhenti di sana, yaitu satu-satunya suhu ketika es
dan air berdampingan pada \qty{1}{atm}: pembekuannya terhenti begitu
kalor yang dilepaskannya membawa campurannya ke sana.

\textbf{22.} $x = c\,\Delta T/L_f = 4.18 \times 8/334 = 0.10$.

\textbf{23.} $\Delta s = c\ln(273/265) - xL_f/273 = 124.3 - 122.3 =
\qty{+2.0}{J/(kg.K)}$: adiabatik, sehingga semuanya tercipta; positif, sebagaimana
seharusnya bagi proses yang spontan.

\textbf{24.} $x = 4.18 \times 5/2257 = 0.93\%$: \qty{9.3}{g} uap tiap liter,
$9.3 \times 10^{-3} \times 1.67 = \qty{16}{L}$ uap yang lahir di dalam badannya dalam
sepersekian sekon: cangkirnya meletup.

\textbf{25.} $T_0S_{\mathrm{cipta}} = 273 \times 2.0 = \qty{550}{J}$ tiap liter,
berbanding $mgh = \qty{9.8}{J}$: lima puluh enam kali lebih besar. \omterm{def:b1:phase-changes:metastable}{Keadaan metastabil}
menyimpan usaha --- yaitu apa yang dapat diambil sebuah \omterm{def:b1:point-kinematics:frame}{lintasan} reversibel --- dan sebuah
ketukan mengubah semuanya menjadi entropi.
\end{solution}
